\section*{Erd\H{o}s Problem 649: smooth numbers preceding a prescribed prime factor}

\paragraph{Statement.}
For primes \(p<q\), must there exist an integer \(n\) whose greatest prime
factor is \(p\), while the greatest prime factor of \(n+1\) is \(q\)?

\paragraph{Claim and attribution.}
No. We give a complete elementary counterexample and then a stronger
family using Dirichlet's theorem and quadratic reciprocity.
The problem and the congruence obstruction are recorded at
\url{https://www.erdosproblems.com/649}. This is a reconstruction of a
known disproof.

\paragraph{An explicit counterexample.}
Take \(p=2\), \(q=7\). An integer with greatest prime factor \(2\)
is \(2^a\), \(a\geq1\). Its residue modulo \(7\) lies in
\(\{1,2,4\}\), since \(2^3\equiv1\pmod7\).
Consequently \(7\nmid 2^a+1\). In particular the greatest prime factor
of \(2^a+1\) cannot be \(7\). This alone disproves the statement.

\paragraph{Infinitely many obstructions for each fixed \(p\).}
Let
\[
 M=8\prod_{\substack{\ell\leq p\\\ell\ {\rm odd\ prime}}}\ell.
\]
Dirichlet's theorem supplies infinitely many primes \(q\equiv-1\pmod M\);
discard the finitely many \(q\leq p\). We claim that every prime
\(\ell\leq p\) is a quadratic residue modulo \(q\).
For \(\ell=2\), this follows from \(q\equiv7\pmod8\).
For odd \(\ell\leq p\), quadratic reciprocity and \(q\equiv3\pmod4\)
give
\[
 \left(\frac{\ell}{q}\right)
 =(-1)^{(\ell-1)/2}\left(\frac{q}{\ell}\right)
 =(-1)^{(\ell-1)/2}\left(\frac{-1}{\ell}\right)=1.
\]
Thus every positive integer all of whose prime factors are at most
\(p\) is a nonzero quadratic residue modulo \(q\).
But \(-1\) is a quadratic nonresidue because \(q\equiv3\pmod4\).
No such integer can satisfy \(n\equiv-1\pmod q\), so \(q\nmid n+1\).
This proves the stronger obstruction even without requiring that \(p\)
actually divide \(n\).

\paragraph{Scope and dependencies.}
The single counterexample uses only modular arithmetic and completely
answers the question. Only the supplementary infinite-family statement
uses the two classical theorems named above.

\paragraph{Completion estimate.}
\textbf{COMPLETION: 100\%}
